\documentclass[aspectratio=169,10pt,english]{beamer}

\input{../../../_preambulo_beamer.tex}
\input{../../../_comandos_beamer.tex}

\selectlanguage{english}

\title{MA1001B - Statistical Modeling for Decision Making}
\subtitle{Lesson 17: Poisson Model for Rare Events}
\author{}
\institute{Tecnol\'ogico de Monterrey}
\date{}

\begin{document}

\begin{frame}[plain]
  \vspace{1.2cm}
  {\Large\bfseries MA1001B - Statistical Modeling for Decision Making\par}
  \vspace{0.7cm}
  {\large Lesson 17: Poisson Model for Rare Events\par}
  \vspace{0.7cm}
  {\color{TecRojo}\rule{\textwidth}{1.2pt}}
  \vspace{0.8cm}

  MA1001B Faculty\\[0.3cm]
  Term\\[0.3cm]
  Tecnol\'ogico de Monterrey
\end{frame}

\begin{frame}{Counts in a Shift}
  \begin{alertblock}{Guiding question}
    If stoppages occur at mean rate $\lambda=3.2$ per shift, what are
    $P(X=0)$ and $P(X\ge 6)$, and what happens if events cluster?
  \end{alertblock}

  \vspace{0.6em}
  By the end of this lesson, you can:
  \begin{itemize}
    \item use the Poisson mean-equals-variance property;
    \item compute $P(X=0)$ and $P(X\ge 6)$;
    \item compare those values with a seed-$42$ simulation;
    \item explain overdispersion when events cluster.
  \end{itemize}

  \vfill
  \scriptsize All shifts used today are synthetic. No real company data are used.
\end{frame}

\begin{frame}{Poisson($\lambda=3.2$)}
  \small
  $X$ is the number of events in one shift.
  \[
    E[X]=\lambda=3.2,\qquad \Var(X)=\lambda=3.2,
  \]
  \[
    P(X=0)=e^{-3.2}=0.040762,\qquad P(X\ge 6)=0.105408.
  \]
\end{frame}

\begin{frame}[fragile]{Step 1: Poisson PMF}
  \begin{columns}[T]
    \begin{column}{0.42\textwidth}
      \small
      \begin{lstlisting}
model = stats.poisson(mu=3.2)
p0 = model.pmf(0)
p6 = model.sf(5)
      \end{lstlisting}

      \begin{block}{Verified output}
        $\lambda=3.200000$\\
        $P(X=0)=0.040762$\\
        $P(X\ge 6)=0.105408$
      \end{block}
    \end{column}
    \begin{column}{0.58\textwidth}
      \centering
      \includegraphics[width=\textwidth]{../figures/poisson_01_pmf.png}
      \vspace{0.2em}
      \scriptsize Poisson PMF from Step 1
    \end{column}
  \end{columns}
\end{frame}

\begin{frame}{Ten Thousand Simulated Shifts}
  \small
  Seed $42$, $n=10{,}000$ Poisson draws:
  \[
    \widehat{P}(X=0)=0.040000,\qquad
    \widehat{P}(X\ge 6)=0.109000.
  \]
  Simulated mean $3.209300$ and variance $3.256094$ sit near $\lambda=3.2$.
\end{frame}

\begin{frame}[fragile]{Step 2: Simulation Check}
  \begin{columns}[T]
    \begin{column}{0.40\textwidth}
      \small
      \begin{lstlisting}
rng = np.random.default_rng(42)
x = rng.poisson(3.2, 10000)
      \end{lstlisting}

      \begin{block}{Verified output}
        Simulated $P(X=0)=0.040000$\\
        Simulated $P(X\ge 6)=0.109000$\\
        Mean $=3.209300$
      \end{block}
    \end{column}
    \begin{column}{0.60\textwidth}
      \centering
      \includegraphics[width=\textwidth]{../figures/poisson_02_simulate_shifts.png}
      \vspace{0.2em}
      \scriptsize Histogram versus PMF from Step 2
    \end{column}
  \end{columns}
\end{frame}

\begin{frame}{Step 3: Clustering Produces Overdispersion}
  \begin{columns}[T]
    \begin{column}{0.42\textwidth}
      \small
      Quiet shifts $\lambda=2.0$ ($70\%$) and busy shifts $\lambda=6.0$
      ($30\%$) keep mean $3.2$ but raise variance to $6.560000$ in theory and
      $6.811382$ in the seed-$42$ mixture.

      \vspace{0.4em}
      \textbf{Limit:} if events cluster, variance exceeds $\lambda$.
      Lesson 18 models waiting times until the first or $r$-th success.
    \end{column}
    \begin{column}{0.58\textwidth}
      \centering
      \includegraphics[width=\textwidth]{../figures/poisson_03_overdispersion.png}
      \vspace{0.2em}
      \scriptsize Overdispersion from Step 3
    \end{column}
  \end{columns}
\end{frame}

\begin{frame}{Concept Check}
  \begin{enumerate}
    \item What is $\mathrm{Var}(X)$ for a Poisson($3.2$) random variable?
    \item Compute $P(X=0)$ from $e^{-3.2}$.
    \item Why is a simulated rate of $0.040000$ not a replacement for $0.040762$?
    \item Why does a mixture of $\lambda=2$ and $\lambda=6$ overdisperse?
  \end{enumerate}

  \vspace{0.6em}
  \begin{block}{Connection to Lesson 18}
    The session can continue directly into geometric and negative-binomial
    waiting times.
  \end{block}

  \vfill
  \scriptsize Source: OpenStax, \textit{Introductory Business Statistics 2e},
  Section 4.4. CC BY-NC-SA.
\end{frame}

\appendix



\end{document}
